%% Приложение А. Описание используемого программного стека
%% Включается в основной документ командой %% Приложение А. Описание используемого программного стека
%% Включается в основной документ командой %% Приложение А. Описание используемого программного стека
%% Включается в основной документ командой %% Приложение А. Описание используемого программного стека
%% Включается в основной документ командой \input{appendix_a_stack}

\appendix
\section*{Приложение А. Описание используемого программного стека}
\label{appendix:stack}
\addcontentsline{toc}{section}{Приложение А. Описание используемого программного стека}

Настоящее приложение содержит сводку инструментов и библиотек, применявшихся
в ходе выполнения работы. Выбор каждого компонента обусловлен функциональными
возможностями, совместимостью с остальным стеком и распространённостью
в сообществе исследователей генеративных моделей. Сводный перечень приведён
в таблице~\ref{tab:stack}; ниже даётся краткое обоснование для компонент,
играющих принципиальную роль в реализации именно метода B-LoRA.
Полная конфигурация проекта (Poetry-зависимости, Hydra-конфиги экспериментов,
DAG-и Airflow, скрипты обучения и оценки) опубликована в открытом репозитории
\url{https://github.com/betyavan/b-lora-flux}.

\begin{table}[h]
    \centering
    \footnotesize
    \caption{Программный стек проекта}
    \label{tab:stack}
    \begin{tabular}{p{3.6cm} p{1.8cm} p{8.0cm}}
        \toprule
        \textbf{Компонент} & \textbf{Версия} & \textbf{Назначение} \\
        \midrule
        \multicolumn{3}{l}{\textit{Машинное обучение}} \\
        Python & 3.10 & Основной язык реализации \\
        PyTorch~\cite{paszke2019pytorch} & 2.1.0 & Вычислительный граф, autograd, GPU \\
        torchvision & 0.16.0 & Трансформации изображений, VGG для LPIPS \\
        Diffusers & $\geq$0.30.0 & \texttt{FluxPipeline}, загрузка LoRA-весов \\
        Transformers~\cite{wolf2020transformers} & $\geq$4.44.0 & Текстовые энкодеры CLIP-L и T5-XXL \\
        Accelerate & $\geq$0.33.0 & Mixed-precision, gradient checkpointing \\
        PEFT~\cite{hu2022lora} & $\geq$0.12.0 & LoRA через \texttt{LoraConfig.target\_modules} \\
        SafeTensors & $\geq$0.4.0 & Безопасное хранение весов адаптеров \\
        SentencePiece & $\geq$0.2.0 & Токенизация для T5-XXL \\
        \midrule
        \multicolumn{3}{l}{\textit{Данные и метрики}} \\
        Datasets & $\geq$2.20.0 & Загрузка ArtBench-10~\cite{liao2022artbench} с HF Hub \\
        Pillow & $\geq$10.0.0 & Предобработка изображений \\
        TorchMetrics & $\geq$1.4.0 & PSNR, SSIM, CLIP-Score \\
        lpips~\cite{zhang2018lpips} & $\geq$0.1.4 & LPIPS на признаках AlexNet/VGG \\
        clean-fid~\cite{parmar2022cleanfid} & $\geq$0.1.35 & Корректный расчёт FID \\
        \midrule
        \multicolumn{3}{l}{\textit{Конфигурация и трекинг}} \\
        Hydra~\cite{von2022hydra} & 1.3.2 & Иерархические YAML-конфиги, multirun \\
        OmegaConf & 2.3.0 & Интерполяция значений в Hydra \\
        ClearML~\cite{clearml2024} & 1.16.4 & Трекинг гиперпараметров, метрик, артефактов \\
        \midrule
        \multicolumn{3}{l}{\textit{Версионирование}} \\
        Git & --- & Контроль версий кода и LaTeX-исходников \\
        DVC~\cite{dvc2020} & 3.54.0 & Версионирование данных и весов через S3 \\
        \midrule
        \multicolumn{3}{l}{\textit{Качество кода}} \\
        Poetry & $\geq$2.0 & Детерминированные зависимости (\texttt{poetry.lock}) \\
        Black & 24.1.1 & Автоформатирование (длина строки 120) \\
        Ruff & 0.6.9 & Линтер на Rust (PEP~8, isort) \\
        Mypy & 1.8.0 & Статическая проверка типов (\texttt{--strict}) \\
        Pytest & 8.3.0 & Юнит- и интеграционные тесты \\
        pre-commit + detect-secrets & 3.0 / 1.5.0 & Защита от утечки секретов в Git \\
        \midrule
        \multicolumn{3}{l}{\textit{Обучение и инфраструктура}} \\
        AI Toolkit (ostris) & 0.9.4 & Бэкенд обучения LoRA для FLUX.1 \\
        Apache Airflow & 2.x & Оркестрация пайплайна train$\to$generate$\to$metrics \\
        Boto3 & $\geq$1.35 & S3-доступ внутри Python-задач Airflow \\
        s3cmd & --- & Синхронизация бинарных артефактов в shell-скриптах \\
        \bottomrule
    \end{tabular}
\end{table}

\subsection*{А.1. AI Toolkit как базовый тренировочный бэкенд}

AI~Toolkit (\texttt{ostris/ai-toolkit})~--- открытый фреймворк обучения
диффузионных моделей с явной поддержкой FLUX.1-dev/schnell. Принципиальная
особенность, обусловившая его выбор для настоящей работы,~--- возможность
указывать целевые модули LoRA на уровне отдельных трансформерных блоков
через YAML-конфигурацию. Это позволяет реализовать поблочное разделение
адаптеров $\Theta_\text{style}$ и $\Theta_\text{content}$ без написания
низкоуровневого кода обучения. Альтернативные решения (kohya-ss/sd-scripts,
diffusers~Trainer) такой гранулярности не предоставляют. Репозиторий
подключён к проекту как git-субмодуль (\texttt{src/ai-toolkit}), что
фиксирует используемую ревизию.

\subsection*{А.2. ClearML и DVC для воспроизводимости}

Платформа ClearML~\cite{clearml2024} обеспечивает автоматическую фиксацию
гиперпараметров, метрик (DINO-style, CLIP-style, CLIP-content, FID, LPIPS),
артефактов (чекпоинты, конфигурации) и системных характеристик для каждого
запуска. Возможность полностью локального развёртывания сервера выгодно
отличает её от Weights~\&~Biases в условиях работы с закрытыми датасетами.

DVC~\cite{dvc2020} расширяет Git версионированием крупных бинарных объектов
(изображений, весов, результатов инференса) с хранением содержимого
в S3-совместимом удалённом хранилище. Механизм пайплайнов
(\texttt{dvc.yaml}) фиксирует ориентированный ациклический граф зависимостей
между этапами и обеспечивает инкрементальное воспроизведение только
изменившихся участков. Команда \texttt{dvc~pull} восстанавливает идентичный
снимок данных после клонирования репозитория.

\subsection*{А.3. Apache Airflow и удалённое исполнение}

Apache~Airflow~2.x управляет последовательностью этапов
train~\textrightarrow{}~generate~\textrightarrow{}~metrics для каждого
эксперимента. Задачи описываются декоратором \texttt{@task}, а отправка
заданий на удалённый узел H100 реализована пользовательскими операторами
\texttt{JobSubmitOperatorEnv} и \texttt{JobSubmitDictOperator}. Динамическое
маппирование (\texttt{.partial().expand()}) обеспечивает параллельный запуск
группы экспериментов без дублирования кода DAG. Передача артефактов между
этапами выполняется через S3: внутри shell-скриптов используется
\texttt{s3cmd~sync} (инкрементальная rsync-подобная семантика, не требует
Python-окружения на узле), а в Python-задачах Airflow~--- \texttt{boto3}
с прямой интеграцией через XCom.


\appendix
\section*{Приложение А. Описание используемого программного стека}
\label{appendix:stack}
\addcontentsline{toc}{section}{Приложение А. Описание используемого программного стека}

Настоящее приложение содержит сводку инструментов и библиотек, применявшихся
в ходе выполнения работы. Выбор каждого компонента обусловлен функциональными
возможностями, совместимостью с остальным стеком и распространённостью
в сообществе исследователей генеративных моделей. Сводный перечень приведён
в таблице~\ref{tab:stack}; ниже даётся краткое обоснование для компонент,
играющих принципиальную роль в реализации именно метода B-LoRA.
Полная конфигурация проекта (Poetry-зависимости, Hydra-конфиги экспериментов,
DAG-и Airflow, скрипты обучения и оценки) опубликована в открытом репозитории
\url{https://github.com/betyavan/b-lora-flux}.

\begin{table}[h]
    \centering
    \footnotesize
    \caption{Программный стек проекта}
    \label{tab:stack}
    \begin{tabular}{p{3.6cm} p{1.8cm} p{8.0cm}}
        \toprule
        \textbf{Компонент} & \textbf{Версия} & \textbf{Назначение} \\
        \midrule
        \multicolumn{3}{l}{\textit{Машинное обучение}} \\
        Python & 3.10 & Основной язык реализации \\
        PyTorch~\cite{paszke2019pytorch} & 2.1.0 & Вычислительный граф, autograd, GPU \\
        torchvision & 0.16.0 & Трансформации изображений, VGG для LPIPS \\
        Diffusers & $\geq$0.30.0 & \texttt{FluxPipeline}, загрузка LoRA-весов \\
        Transformers~\cite{wolf2020transformers} & $\geq$4.44.0 & Текстовые энкодеры CLIP-L и T5-XXL \\
        Accelerate & $\geq$0.33.0 & Mixed-precision, gradient checkpointing \\
        PEFT~\cite{hu2022lora} & $\geq$0.12.0 & LoRA через \texttt{LoraConfig.target\_modules} \\
        SafeTensors & $\geq$0.4.0 & Безопасное хранение весов адаптеров \\
        SentencePiece & $\geq$0.2.0 & Токенизация для T5-XXL \\
        \midrule
        \multicolumn{3}{l}{\textit{Данные и метрики}} \\
        Datasets & $\geq$2.20.0 & Загрузка ArtBench-10~\cite{liao2022artbench} с HF Hub \\
        Pillow & $\geq$10.0.0 & Предобработка изображений \\
        TorchMetrics & $\geq$1.4.0 & PSNR, SSIM, CLIP-Score \\
        lpips~\cite{zhang2018lpips} & $\geq$0.1.4 & LPIPS на признаках AlexNet/VGG \\
        clean-fid~\cite{parmar2022cleanfid} & $\geq$0.1.35 & Корректный расчёт FID \\
        \midrule
        \multicolumn{3}{l}{\textit{Конфигурация и трекинг}} \\
        Hydra~\cite{von2022hydra} & 1.3.2 & Иерархические YAML-конфиги, multirun \\
        OmegaConf & 2.3.0 & Интерполяция значений в Hydra \\
        ClearML~\cite{clearml2024} & 1.16.4 & Трекинг гиперпараметров, метрик, артефактов \\
        \midrule
        \multicolumn{3}{l}{\textit{Версионирование}} \\
        Git & --- & Контроль версий кода и LaTeX-исходников \\
        DVC~\cite{dvc2020} & 3.54.0 & Версионирование данных и весов через S3 \\
        \midrule
        \multicolumn{3}{l}{\textit{Качество кода}} \\
        Poetry & $\geq$2.0 & Детерминированные зависимости (\texttt{poetry.lock}) \\
        Black & 24.1.1 & Автоформатирование (длина строки 120) \\
        Ruff & 0.6.9 & Линтер на Rust (PEP~8, isort) \\
        Mypy & 1.8.0 & Статическая проверка типов (\texttt{--strict}) \\
        Pytest & 8.3.0 & Юнит- и интеграционные тесты \\
        pre-commit + detect-secrets & 3.0 / 1.5.0 & Защита от утечки секретов в Git \\
        \midrule
        \multicolumn{3}{l}{\textit{Обучение и инфраструктура}} \\
        AI Toolkit (ostris) & 0.9.4 & Бэкенд обучения LoRA для FLUX.1 \\
        Apache Airflow & 2.x & Оркестрация пайплайна train$\to$generate$\to$metrics \\
        Boto3 & $\geq$1.35 & S3-доступ внутри Python-задач Airflow \\
        s3cmd & --- & Синхронизация бинарных артефактов в shell-скриптах \\
        \bottomrule
    \end{tabular}
\end{table}

\subsection*{А.1. AI Toolkit как базовый тренировочный бэкенд}

AI~Toolkit (\texttt{ostris/ai-toolkit})~--- открытый фреймворк обучения
диффузионных моделей с явной поддержкой FLUX.1-dev/schnell. Принципиальная
особенность, обусловившая его выбор для настоящей работы,~--- возможность
указывать целевые модули LoRA на уровне отдельных трансформерных блоков
через YAML-конфигурацию. Это позволяет реализовать поблочное разделение
адаптеров $\Theta_\text{style}$ и $\Theta_\text{content}$ без написания
низкоуровневого кода обучения. Альтернативные решения (kohya-ss/sd-scripts,
diffusers~Trainer) такой гранулярности не предоставляют. Репозиторий
подключён к проекту как git-субмодуль (\texttt{src/ai-toolkit}), что
фиксирует используемую ревизию.

\subsection*{А.2. ClearML и DVC для воспроизводимости}

Платформа ClearML~\cite{clearml2024} обеспечивает автоматическую фиксацию
гиперпараметров, метрик (DINO-style, CLIP-style, CLIP-content, FID, LPIPS),
артефактов (чекпоинты, конфигурации) и системных характеристик для каждого
запуска. Возможность полностью локального развёртывания сервера выгодно
отличает её от Weights~\&~Biases в условиях работы с закрытыми датасетами.

DVC~\cite{dvc2020} расширяет Git версионированием крупных бинарных объектов
(изображений, весов, результатов инференса) с хранением содержимого
в S3-совместимом удалённом хранилище. Механизм пайплайнов
(\texttt{dvc.yaml}) фиксирует ориентированный ациклический граф зависимостей
между этапами и обеспечивает инкрементальное воспроизведение только
изменившихся участков. Команда \texttt{dvc~pull} восстанавливает идентичный
снимок данных после клонирования репозитория.

\subsection*{А.3. Apache Airflow и удалённое исполнение}

Apache~Airflow~2.x управляет последовательностью этапов
train~\textrightarrow{}~generate~\textrightarrow{}~metrics для каждого
эксперимента. Задачи описываются декоратором \texttt{@task}, а отправка
заданий на удалённый узел H100 реализована пользовательскими операторами
\texttt{JobSubmitOperatorEnv} и \texttt{JobSubmitDictOperator}. Динамическое
маппирование (\texttt{.partial().expand()}) обеспечивает параллельный запуск
группы экспериментов без дублирования кода DAG. Передача артефактов между
этапами выполняется через S3: внутри shell-скриптов используется
\texttt{s3cmd~sync} (инкрементальная rsync-подобная семантика, не требует
Python-окружения на узле), а в Python-задачах Airflow~--- \texttt{boto3}
с прямой интеграцией через XCom.


\appendix
\section*{Приложение А. Описание используемого программного стека}
\label{appendix:stack}
\addcontentsline{toc}{section}{Приложение А. Описание используемого программного стека}

Настоящее приложение содержит сводку инструментов и библиотек, применявшихся
в ходе выполнения работы. Выбор каждого компонента обусловлен функциональными
возможностями, совместимостью с остальным стеком и распространённостью
в сообществе исследователей генеративных моделей. Сводный перечень приведён
в таблице~\ref{tab:stack}; ниже даётся краткое обоснование для компонент,
играющих принципиальную роль в реализации именно метода B-LoRA.
Полная конфигурация проекта (Poetry-зависимости, Hydra-конфиги экспериментов,
DAG-и Airflow, скрипты обучения и оценки) опубликована в открытом репозитории
\url{https://github.com/betyavan/b-lora-flux}.

\begin{table}[h]
    \centering
    \footnotesize
    \caption{Программный стек проекта}
    \label{tab:stack}
    \begin{tabular}{p{3.6cm} p{1.8cm} p{8.0cm}}
        \toprule
        \textbf{Компонент} & \textbf{Версия} & \textbf{Назначение} \\
        \midrule
        \multicolumn{3}{l}{\textit{Машинное обучение}} \\
        Python & 3.10 & Основной язык реализации \\
        PyTorch~\cite{paszke2019pytorch} & 2.1.0 & Вычислительный граф, autograd, GPU \\
        torchvision & 0.16.0 & Трансформации изображений, VGG для LPIPS \\
        Diffusers & $\geq$0.30.0 & \texttt{FluxPipeline}, загрузка LoRA-весов \\
        Transformers~\cite{wolf2020transformers} & $\geq$4.44.0 & Текстовые энкодеры CLIP-L и T5-XXL \\
        Accelerate & $\geq$0.33.0 & Mixed-precision, gradient checkpointing \\
        PEFT~\cite{hu2022lora} & $\geq$0.12.0 & LoRA через \texttt{LoraConfig.target\_modules} \\
        SafeTensors & $\geq$0.4.0 & Безопасное хранение весов адаптеров \\
        SentencePiece & $\geq$0.2.0 & Токенизация для T5-XXL \\
        \midrule
        \multicolumn{3}{l}{\textit{Данные и метрики}} \\
        Datasets & $\geq$2.20.0 & Загрузка ArtBench-10~\cite{liao2022artbench} с HF Hub \\
        Pillow & $\geq$10.0.0 & Предобработка изображений \\
        TorchMetrics & $\geq$1.4.0 & PSNR, SSIM, CLIP-Score \\
        lpips~\cite{zhang2018lpips} & $\geq$0.1.4 & LPIPS на признаках AlexNet/VGG \\
        clean-fid~\cite{parmar2022cleanfid} & $\geq$0.1.35 & Корректный расчёт FID \\
        \midrule
        \multicolumn{3}{l}{\textit{Конфигурация и трекинг}} \\
        Hydra~\cite{von2022hydra} & 1.3.2 & Иерархические YAML-конфиги, multirun \\
        OmegaConf & 2.3.0 & Интерполяция значений в Hydra \\
        ClearML~\cite{clearml2024} & 1.16.4 & Трекинг гиперпараметров, метрик, артефактов \\
        \midrule
        \multicolumn{3}{l}{\textit{Версионирование}} \\
        Git & --- & Контроль версий кода и LaTeX-исходников \\
        DVC~\cite{dvc2020} & 3.54.0 & Версионирование данных и весов через S3 \\
        \midrule
        \multicolumn{3}{l}{\textit{Качество кода}} \\
        Poetry & $\geq$2.0 & Детерминированные зависимости (\texttt{poetry.lock}) \\
        Black & 24.1.1 & Автоформатирование (длина строки 120) \\
        Ruff & 0.6.9 & Линтер на Rust (PEP~8, isort) \\
        Mypy & 1.8.0 & Статическая проверка типов (\texttt{--strict}) \\
        Pytest & 8.3.0 & Юнит- и интеграционные тесты \\
        pre-commit + detect-secrets & 3.0 / 1.5.0 & Защита от утечки секретов в Git \\
        \midrule
        \multicolumn{3}{l}{\textit{Обучение и инфраструктура}} \\
        AI Toolkit (ostris) & 0.9.4 & Бэкенд обучения LoRA для FLUX.1 \\
        Apache Airflow & 2.x & Оркестрация пайплайна train$\to$generate$\to$metrics \\
        Boto3 & $\geq$1.35 & S3-доступ внутри Python-задач Airflow \\
        s3cmd & --- & Синхронизация бинарных артефактов в shell-скриптах \\
        \bottomrule
    \end{tabular}
\end{table}

\subsection*{А.1. AI Toolkit как базовый тренировочный бэкенд}

AI~Toolkit (\texttt{ostris/ai-toolkit})~--- открытый фреймворк обучения
диффузионных моделей с явной поддержкой FLUX.1-dev/schnell. Принципиальная
особенность, обусловившая его выбор для настоящей работы,~--- возможность
указывать целевые модули LoRA на уровне отдельных трансформерных блоков
через YAML-конфигурацию. Это позволяет реализовать поблочное разделение
адаптеров $\Theta_\text{style}$ и $\Theta_\text{content}$ без написания
низкоуровневого кода обучения. Альтернативные решения (kohya-ss/sd-scripts,
diffusers~Trainer) такой гранулярности не предоставляют. Репозиторий
подключён к проекту как git-субмодуль (\texttt{src/ai-toolkit}), что
фиксирует используемую ревизию.

\subsection*{А.2. ClearML и DVC для воспроизводимости}

Платформа ClearML~\cite{clearml2024} обеспечивает автоматическую фиксацию
гиперпараметров, метрик (DINO-style, CLIP-style, CLIP-content, FID, LPIPS),
артефактов (чекпоинты, конфигурации) и системных характеристик для каждого
запуска. Возможность полностью локального развёртывания сервера выгодно
отличает её от Weights~\&~Biases в условиях работы с закрытыми датасетами.

DVC~\cite{dvc2020} расширяет Git версионированием крупных бинарных объектов
(изображений, весов, результатов инференса) с хранением содержимого
в S3-совместимом удалённом хранилище. Механизм пайплайнов
(\texttt{dvc.yaml}) фиксирует ориентированный ациклический граф зависимостей
между этапами и обеспечивает инкрементальное воспроизведение только
изменившихся участков. Команда \texttt{dvc~pull} восстанавливает идентичный
снимок данных после клонирования репозитория.

\subsection*{А.3. Apache Airflow и удалённое исполнение}

Apache~Airflow~2.x управляет последовательностью этапов
train~\textrightarrow{}~generate~\textrightarrow{}~metrics для каждого
эксперимента. Задачи описываются декоратором \texttt{@task}, а отправка
заданий на удалённый узел H100 реализована пользовательскими операторами
\texttt{JobSubmitOperatorEnv} и \texttt{JobSubmitDictOperator}. Динамическое
маппирование (\texttt{.partial().expand()}) обеспечивает параллельный запуск
группы экспериментов без дублирования кода DAG. Передача артефактов между
этапами выполняется через S3: внутри shell-скриптов используется
\texttt{s3cmd~sync} (инкрементальная rsync-подобная семантика, не требует
Python-окружения на узле), а в Python-задачах Airflow~--- \texttt{boto3}
с прямой интеграцией через XCom.


\appendix
\section*{Приложение А. Описание используемого программного стека}
\label{appendix:stack}
\addcontentsline{toc}{section}{Приложение А. Описание используемого программного стека}

Настоящее приложение содержит сводку инструментов и библиотек, применявшихся
в ходе выполнения работы. Выбор каждого компонента обусловлен функциональными
возможностями, совместимостью с остальным стеком и распространённостью
в сообществе исследователей генеративных моделей. Сводный перечень приведён
в таблице~\ref{tab:stack}; ниже даётся краткое обоснование для компонент,
играющих принципиальную роль в реализации именно метода B-LoRA.
Полная конфигурация проекта (Poetry-зависимости, Hydra-конфиги экспериментов,
DAG-и Airflow, скрипты обучения и оценки) опубликована в открытом репозитории
\url{https://github.com/betyavan/b-lora-flux}.

\begin{table}[h]
    \centering
    \footnotesize
    \caption{Программный стек проекта}
    \label{tab:stack}
    \begin{tabular}{p{3.6cm} p{1.8cm} p{8.0cm}}
        \toprule
        \textbf{Компонент} & \textbf{Версия} & \textbf{Назначение} \\
        \midrule
        \multicolumn{3}{l}{\textit{Машинное обучение}} \\
        Python & 3.10 & Основной язык реализации \\
        PyTorch~\cite{paszke2019pytorch} & 2.1.0 & Вычислительный граф, autograd, GPU \\
        torchvision & 0.16.0 & Трансформации изображений, VGG для LPIPS \\
        Diffusers & $\geq$0.30.0 & \texttt{FluxPipeline}, загрузка LoRA-весов \\
        Transformers~\cite{wolf2020transformers} & $\geq$4.44.0 & Текстовые энкодеры CLIP-L и T5-XXL \\
        Accelerate & $\geq$0.33.0 & Mixed-precision, gradient checkpointing \\
        PEFT~\cite{hu2022lora} & $\geq$0.12.0 & LoRA через \texttt{LoraConfig.target\_modules} \\
        SafeTensors & $\geq$0.4.0 & Безопасное хранение весов адаптеров \\
        SentencePiece & $\geq$0.2.0 & Токенизация для T5-XXL \\
        \midrule
        \multicolumn{3}{l}{\textit{Данные и метрики}} \\
        Datasets & $\geq$2.20.0 & Загрузка ArtBench-10~\cite{liao2022artbench} с HF Hub \\
        Pillow & $\geq$10.0.0 & Предобработка изображений \\
        TorchMetrics & $\geq$1.4.0 & PSNR, SSIM, CLIP-Score \\
        lpips~\cite{zhang2018lpips} & $\geq$0.1.4 & LPIPS на признаках AlexNet/VGG \\
        clean-fid~\cite{parmar2022cleanfid} & $\geq$0.1.35 & Корректный расчёт FID \\
        \midrule
        \multicolumn{3}{l}{\textit{Конфигурация и трекинг}} \\
        Hydra~\cite{von2022hydra} & 1.3.2 & Иерархические YAML-конфиги, multirun \\
        OmegaConf & 2.3.0 & Интерполяция значений в Hydra \\
        ClearML~\cite{clearml2024} & 1.16.4 & Трекинг гиперпараметров, метрик, артефактов \\
        \midrule
        \multicolumn{3}{l}{\textit{Версионирование}} \\
        Git & --- & Контроль версий кода и LaTeX-исходников \\
        DVC~\cite{dvc2020} & 3.54.0 & Версионирование данных и весов через S3 \\
        \midrule
        \multicolumn{3}{l}{\textit{Качество кода}} \\
        Poetry & $\geq$2.0 & Детерминированные зависимости (\texttt{poetry.lock}) \\
        Black & 24.1.1 & Автоформатирование (длина строки 120) \\
        Ruff & 0.6.9 & Линтер на Rust (PEP~8, isort) \\
        Mypy & 1.8.0 & Статическая проверка типов (\texttt{--strict}) \\
        Pytest & 8.3.0 & Юнит- и интеграционные тесты \\
        pre-commit + detect-secrets & 3.0 / 1.5.0 & Защита от утечки секретов в Git \\
        \midrule
        \multicolumn{3}{l}{\textit{Обучение и инфраструктура}} \\
        AI Toolkit (ostris) & 0.9.4 & Бэкенд обучения LoRA для FLUX.1 \\
        Apache Airflow & 2.x & Оркестрация пайплайна train$\to$generate$\to$metrics \\
        Boto3 & $\geq$1.35 & S3-доступ внутри Python-задач Airflow \\
        s3cmd & --- & Синхронизация бинарных артефактов в shell-скриптах \\
        \bottomrule
    \end{tabular}
\end{table}

\subsection*{А.1. AI Toolkit как базовый тренировочный бэкенд}

AI~Toolkit (\texttt{ostris/ai-toolkit})~--- открытый фреймворк обучения
диффузионных моделей с явной поддержкой FLUX.1-dev/schnell. Принципиальная
особенность, обусловившая его выбор для настоящей работы,~--- возможность
указывать целевые модули LoRA на уровне отдельных трансформерных блоков
через YAML-конфигурацию. Это позволяет реализовать поблочное разделение
адаптеров $\Theta_\text{style}$ и $\Theta_\text{content}$ без написания
низкоуровневого кода обучения. Альтернативные решения (kohya-ss/sd-scripts,
diffusers~Trainer) такой гранулярности не предоставляют. Репозиторий
подключён к проекту как git-субмодуль (\texttt{src/ai-toolkit}), что
фиксирует используемую ревизию.

\subsection*{А.2. ClearML и DVC для воспроизводимости}

Платформа ClearML~\cite{clearml2024} обеспечивает автоматическую фиксацию
гиперпараметров, метрик (DINO-style, CLIP-style, CLIP-content, FID, LPIPS),
артефактов (чекпоинты, конфигурации) и системных характеристик для каждого
запуска. Возможность полностью локального развёртывания сервера выгодно
отличает её от Weights~\&~Biases в условиях работы с закрытыми датасетами.

DVC~\cite{dvc2020} расширяет Git версионированием крупных бинарных объектов
(изображений, весов, результатов инференса) с хранением содержимого
в S3-совместимом удалённом хранилище. Механизм пайплайнов
(\texttt{dvc.yaml}) фиксирует ориентированный ациклический граф зависимостей
между этапами и обеспечивает инкрементальное воспроизведение только
изменившихся участков. Команда \texttt{dvc~pull} восстанавливает идентичный
снимок данных после клонирования репозитория.

\subsection*{А.3. Apache Airflow и удалённое исполнение}

Apache~Airflow~2.x управляет последовательностью этапов
train~\textrightarrow{}~generate~\textrightarrow{}~metrics для каждого
эксперимента. Задачи описываются декоратором \texttt{@task}, а отправка
заданий на удалённый узел H100 реализована пользовательскими операторами
\texttt{JobSubmitOperatorEnv} и \texttt{JobSubmitDictOperator}. Динамическое
маппирование (\texttt{.partial().expand()}) обеспечивает параллельный запуск
группы экспериментов без дублирования кода DAG. Передача артефактов между
этапами выполняется через S3: внутри shell-скриптов используется
\texttt{s3cmd~sync} (инкрементальная rsync-подобная семантика, не требует
Python-окружения на узле), а в Python-задачах Airflow~--- \texttt{boto3}
с прямой интеграцией через XCom.
